\section*{Abstract}
Drug discovery is a prolonged and expensive procedure, with the identification of viable drug-target interactions (DTIs) being one of its most resource-intensive stages.
Computational screening offers a scalable alternative to first-stage wet-lab screening.
Recent advances in deep learning have improved the accuracy of DTI models trained on structural and sequence-based data.
Many existing models, however, use complex architectures that limit reproducibility and interpretability for users.

This work presents the full pipeline of DTI prediction designed around accessibility and competitive performance.
Drug molecules are encoded as MACCS fingerprints and target proteins as ProtBERT embeddings, with interactions and raw data extracted from the BindingDB dataset (908{,}643 high-confidence pairs).
A subset of interactions classified as weak binders/non-binders, $5.3 \leq \mathrm{pAffinity} \leq 7.0$ defined as the \textit{grey area}, is removed from the dataset to reduce label noise.
A fully connected deep neural network is trained in PyTorch using the Adam optimiser and a weighted binary cross-entropy loss to address class imbalance, with a 65:20:15 train/validation/test split.

The model achieves an AUC-ROC of approximately 0.97 and an F1 score of 0.93 at decision threshold $\tau = 0.3$, exceeding the headline AUCs reported by MINDG (0.95) and DeepCDA (0.89), though direct comparison is limited by differences in dataset scale and evaluation protocol.
The model is served through a FastAPI backend and a React frontend as a virtual screening tool, allowing users to query drug molecules against selected protein targets; the deployed interface screens 100 compounds against a target in under 65\,ms on CPU.
This work demonstrates that a well-tuned fully connected architecture, paired with established molecular and protein encoders, can achieve competitive performance without the overhead of graph- or attention-based designs.
It also highlights the methodological care required when comparing results across evaluation protocols.
\section*{Acknowledgements}
I would like to thank my supervisor, Dr Mohamad Saada, for his guidance and support during the project, and my friends and family for the encouragement throughout my university career.